\normalsize\large\newpage
\chapter{Cation exchange}
\def\shorttitle{Ion exchange} % Abbreviated chapter title for right top header
\section{Background}

Sorption of species on the surface of solids is an important regulating mechanism for concentrations of dissolved ions in natural waters. Natural substances whose surfaces can act as sorbers include clay minerals, organic particles and oxides/hydroxides. The capability of reactive transport models to simulate sorption processes is essential to the successful application of these models.

When dealing with sorption processes a distinction is often made between surfaces with a constant exchange capacity (ion exchange) and surfaces with a variable charge (surface complexation). In ion exchange problems, ions are adsorbed and released in equivalent proportions. The exchange capacity of the exchanging surface is assumed constant and the net charge of the surface does not change during the exchange process. This approach works especially well for cation exchange on clay and organic surfaces. In surface complexation, on the other hand, the charge of the surface is variable and dependent on the amount and kind of ions sorbed. It applies for example to sorption of heavy metals on the surface of oxides and hydroxides. 

PHREEQC can simulate sorption through ion exchange and surface complexation and both approaches will be discussed below.

Cations readily adsorb to negatively charged clay and organic particles. A change of the composition of the water that is in contact with the exchange complex induces exchange and hence a change in the composition of both the water and the exchange complex. Since the amounts of adsorbed ions can be large compared to the amounts in solution, the effect of this process on the water composition can be significant.

Cation exchange occurs in a variety of natural systems but is especially relevant in coastal areas. In seawater Na, Mg and K are present at increased levels compared to fresh water, which generally contains Ca as the dominant cation. When the fresh water displaces the seawater Na, Mg and K from the sediment's exchange complex are displaced by Ca. The exchange of \na\ (on the exchange complex) for \ca\ that occurs in a freshening aquifer is described by the following equilibrium reaction:
\begin{equation}
\half \ca + \nax \leftrightarrow \half \cax + \na
\end{equation}
where X indicates the soil exchanger.

As the affinity of the exchanger differs for each individual cation, a spatial separation of the displaced cations can develop. Going upstream of the fresh- saltwater interface a so-called chromatographic sequence of Na, K, Mg and Ca dominated water types is found (Appelo and Postma, 1994). 

The reverse occurs when seawater intrudes in a fresh water aquifer:
\begin{equation}
\na + \half \cax \leftrightarrow \nax + \half \ca
\end{equation}

The equilibrium constant for this reaction is given by:
\begin{equation}
K_{\text{Na} \backslash \text{Ca}} = \frac{[\beta_{\text{Na}}][\ca]^{\half}}{[\beta_{\text{Ca}}]^{\half}[\na]} = 0.398
\label{eq:knaca}
\end{equation}
where $[\beta]$ indicates the activity of the sorbed species and the subscripts Na and Ca denote the sorbed species $\nax$ and $\cax$, respectively. The equilibrium constant for cation exchange reactions is also called the selectivity coefficient. Values of selectivity coefficients for cations commonly found in natural systems can be found in Appelo and Postma (2005). They are also included in the PHREEQC-2 database.

PHREEQC uses the Gaines-Thomas convention, which means that the activity of the sorbed species $\beta$ is defined as the equivalent fraction. For example for $\cax$:
\begin{equation}
\beta_{Ca} = \frac{\text{meq}_{\cax}}{\text{CEC}}
\end{equation}
in which $\text{meq}_{\cax}$ is the adsorbed $\ca$ concentration expressed in meq/100g and CEC is the \textit{cation exchange capacity}, also expressed in meq/100g. The CEC is a measure of the capability of a material to sorb cations and equals the equivalent sum of all sorbed species. Therefore:
\begin{equation}
\sum \beta = 1
\label{eq:betas}
\end{equation}

Note that equilibrium constants (e.g., Eqn. \ref{eq:knaca}) describing ion exchange in the Gaines-Thomas convention use equivalent fraction = 1 for the standard state for each sorbed species but may use either 1 M or 1 m as the standard state for aqueous species.  When using equilibrium constants taken from the literature to describe ion exchange reactions, it is important to make sure they are consistent with these standard states.

Equation \ref{eq:betas} can be conveniently used to calculate the activities of the adsorbed species at a known water composition by combining it with the expressions of the selectivity coefficients (equation \ref{eq:knaca} and similar expressions for other ions). The result is a quadratic expression that can be solved using the abc-formula (Appelo and Postma, 2005).

\section{PHREEQC Exercise 4: Simulating cation exchange reactions}

\subsection{Calculation of the exchange complex composition}
In this exercise you will first calculate the composition of the exchanger complex by hand and then with PHREEQC. For the hand calculations you are provided with the concentration of the major cations in seawater (in mmol/kgw):

\begin{table}[hb]
\centering
\footnotesize
\begin{tabular}{llll}
\toprule
\na	& \ka & \ca & \mg \\
\midrule
485 & 10.6 & 10.7 & 55.1 \\
\bottomrule
\end{tabular}
\end{table}

and the selectivity coefficients of $\ka$ and $\mg$:

\begin{equation}
K_{\text{Na} \backslash \text{K}} = \frac{[\beta_{\text{Na}}][\ka]}{[\beta_{\text{K}}][\na]} = 0.20
\label{eq:knak}
\end{equation}

\begin{equation}
K_{\text{Na} \backslash \text{Mg}} = \frac{[\beta_{\text{Na}}][\mg]^{\half}}{[\beta_{\text{Mg}}]^{\half}[\na]} = 0.501
\label{eq:knamg}
\end{equation}

Rewrite equations \ref{eq:knaca}, \ref{eq:knak} and \ref{eq:knamg} to express $\beta_\text{Ca}$, $\beta_\text{K}$ and $\beta_\text{Mg}$ as functions of $\beta_\text{Na}$ and the solute concentrations. Insert into equation \ref{eq:betas} and solve for $\beta_\text{Na}$. The other values are found by back-substitution. Fill in the values below:

\begin{table}[hb]
\centering
\footnotesize
\begin{tabular}{lllll}
\toprule
$\beta_\text{Na}$	& $\beta_\text{K}$ & $\beta_\text{Ca}$ & $\beta_\text{Mg}$ & $\sum \beta$\\
\midrule
\ldots\ldots & \ldots\ldots & \ldots\ldots & \ldots\ldots & \ldots\ldots \\
\bottomrule
\end{tabular}
\end{table}

You can also use PHREEQC to get the same answer. Enter the cation concentrations as the solution composition and use the keyword \textbf{EXCHANGE} to have PHREEQC calculate the exchange complex. The syntax is as follows:

\begin{quote}
EXCHANGE 1 \\
	X 0.06 \\
	-equilibrate 1 \\
END
\end{quote}

Under \textbf{EXCHANGE}, the concentration of element X is defined. 
This is the number of exchanger sites in moles. 
In this case, there are 60 millimoles of exchanger sites. The number in milliequivalents is the same because the exchanger is assumed to have a single negative charge in PHREEQC. Here, the exchanger is in contact with a solution that has 1 kg of water (the default amount in PHREEQC), or 1 l. So, the cation exchange capacity expressed in meq/l is CEC = 60 meq/l. Can you calculate how much that is in meq/100g? The answer is CEC = \ldots\ldots meq/100g.

The identifier \textbf{-equilibrate} is used to have the exchanger composition calculated by PHREEQC according to equilibrium with a solution. The number after the identifier indicates a solution number, 1 in this case. Alternatively, the solution composition may be entered explicitly by typing the concentrations of the sorbed species. For the purpose of this exercise however, the \textbf{-equilibrate} option is required.

The activities of the sorbed species ($\beta$) are listed in the PHREEQC output file under the heading `Beginning of initial exchange-composition calculations'. Look for the numbers in the column `Equivalent fractions', which are the $\beta$ values, and insert the appropriate values in the table below.

\begin{table}[h]
\centering
\footnotesize
\begin{tabular}{lllll}
\toprule
 & \na	& \ka & \ca & \mg \\
\midrule
$\beta$ & \ldots\ldots & \ldots\ldots & \ldots\ldots & \ldots\ldots \\
meq/l   & \ldots\ldots & \ldots\ldots & \ldots\ldots & \ldots\ldots \\
$K_d$   & \ldots\ldots & \ldots\ldots & \ldots\ldots & \ldots\ldots \\
\bottomrule
\end{tabular}
\end{table}

Also insert the sorbed concentrations expressed as meq/l of pore water and calculate the distribution coefficient, $K_d$. The distribution coefficient is the ratio of the sorbed concentration over the dissolved concentration and is often used in solute transport studies to estimate a retardation factor. The $K_d$ values for seawater are clearly different from the values for a typical fresh water, listed below, which demonstrates the limited applicability of this approach in settings with a strong concentration gradient.

\begin{table}[h]
\centering
\footnotesize
\begin{tabular}{llll}
\toprule
 & \na	& \ca & \mg \\
\midrule
$\beta$           & 0.0169 & 0.9291 & 0.054 \\
sorbed (meq/l)    & 1.014  & 55.75  & 3.24 \\
solute (mmol/kgw) & 2.43   & 3.04   & 0.28 \\
$K_d$             & 0.417  & 9.17   & 5.79 \\
\bottomrule
\end{tabular}
\end{table}

As a final part of this exercise you can also examine the effect of diluting 
a solution that is in equilibrium with the exchange complex. 
Because the amount of cations on the exchanger is relatively 
large compared to the amount of dissolved ions in a dilute solution, 
the exchanger tends to work as a buffer. 
Dilution of the solution will be accompanied by an increase of the 
monovalent ions relative to the divalent ions, which can be seen by rewriting equation \ref{eq:knaca}:
\begin{equation}
K_{\text{Na} \backslash \text{Ca}} \frac{[\nax]}{[\cax]^{\half}} = \frac{[\na]}{[\ca]^{\half}}
\end{equation}

The left-hand side is essentially constant when the sorbed concentrations 
are large compared to the solute concentration, 
which means that a 10-fold dilution of \na\ is accompanied by a 100-fold dilution of \ca.

Simulate a 1000-fold dilution of the solution in equilibrium with the 
exchanger from the previous PHREEQC simulation. 
You can do this by appending a new simulation in which you equilibrate 
the exchanger with a new solution that has concentrations 1000 times 
less than the seawater. Include the following line:

\begin{quote}
USE exchange 1
\end{quote}

to equilibrate the new solution with the exchange complex from the previous simulation. 
PHREEQC remembers the composition of the exchange complex and the keyword \textbf{USE} 
ensures that the exchanger will be used again in the new simulation.

Inspect your output file and fill in the table below.

\begin{table}[h]
\centering
\footnotesize
\begin{tabular}{llllll}
\toprule
 & \na	& \ka & \ca & \mg & $\frac{\na}{\sqrt{\ca}}$ \\
\midrule
seawater 				& 485   & 10.6   & 10.7   & 55.1   & 4.69 \\
dill. seawater  & 0.485 & 0.0106 & 0.0107 & 0.0551 & 0.148 \\
dill. + equil.  & \ldots\ldots & \ldots\ldots & \ldots\ldots & \ldots\ldots & \ldots\ldots \\
sorbed (meq/l)  & \ldots\ldots & \ldots\ldots & \ldots\ldots & \ldots\ldots & \\
sorbed (meq/l)  & \ldots\ldots & \ldots\ldots & \ldots\ldots & \ldots\ldots & \\ 
\bottomrule
\end{tabular}
\end{table}

Notice that the ratio $\frac{\na}{\sqrt{\ca}}$ has hardly changed after the diluted solution has equilibrated, just as the composition of the exchange complex.

\subsection{Cation exchange reactions at Cape Cod}

This set of exercises is based on several field-scale reactive transport experiments 
conducted in an anaerobic, Fe-reducing aquifer in 
western Cape Cod, Massachusetts, United States 
\citep{fox_2010,Smith_2017} 
%Jamieson et al., in preparation).  
The major cations are $Na$, $K$, $Mg$, $Ca$ and $Fe(II)$; there are minor 
concentrations of $Mn$, $Sr$, $Ba$, and $ammonium$.  
One of the field experiments involved injection native groundwater amended with $Cs$.  
We will examine equilibration of the groundwater with an exchanger
initially in the $Na$ form.  
We will also examine two different cation exchange models.

Initially we equilibrate the exchanger occupancy with the measured groundwater composition.

\begin{itemize}
\item Create a PHREEQC file with the composition of the native groundwater (Table \ref{tab:capecod}).
\end{itemize}
 
\begin{table}[h]
\begin{center}
\centering
\caption{Groundwater composition for Cape Cod}
\begin{tabular}{clll}
\toprule
Parameter	  & Groundwater	& Cs-spike	& Units \\
Temperature	&     12	    &   12	    & $^\circ$C   \\
      pH	  &    6.6	    &   6.6     & -           \\	
      pe	  &    4.5	    &   4.5	    & -           \\
     C(4)	  &   1103	    &  1103	    & $\mu$mol/L  \\
      Na	  &    580	    &   580	    & $\mu$mol/L  \\
K	  & 75	& 75	& $\mu$mol/L \\
Mg	& 53	& 53	& $\mu$mol/L \\
Ca	& 170 &	170	& $\mu$mol/L \\
Fe(2)	& 175	& 175	& $\mu$mol/L  \\
Mn(2)	& 2.2	& 2.2	& $\mu$mol/L  \\
Amm	  & 1.1	& 1.1 &	$\mu$mol/L  \\
Sr	  & 0.08	& 0.08	& $\mu$mol/L  \\
Ba	  & 0.03	& 0.03	& $\mu$mol/L  \\
Cl	  & 482.8	& 482.8	& $\mu$mol/L  \\
S(6)	& 100	  & 100	  & $\mu$mol/L  \\
P	    & 63	  & 63	  & $\mu$mol/L  \\
Cs	  & 0	    & 10	  & $\mu$mol/L  \\ 
\bottomrule
\end{tabular}
\label{tab:capecod}
\end{center}
\end{table}

Before running the simulation, you will need to add thermodynamic information 
for Cs because it is not in the database you will be using (Amm.dat).  
To add an element to a simulation that is not in the database, you need to use the keywords:

\begin{quote}
{\color{red}SOLUTION\_MASTER\_SPECIES \\
\# 1 2 3 4 5 \\
  Cs	Cs+	0	Cs	132.905 \\
\#1: Constituent name \\
\#2: Master species for the elements \\
\#3: Contribution to the alkalinity \\
\#4: Formula unit for the formula weight \\
\#5: gram formula weight for the formula unit \\
SOLUTION\_SPECIES \\
Cs+  =  Cs+ \\
log\_k  0 }\\
\end{quote}

The entry under solution species is needed to build the master species into the database.  
The logK of 0 is the default so could be left out.

Once that information is added, the simulation can be run using the database Amm.dat. 
The Amm.dat database is identical to the Phreeqc.dat
database except that $Amm$ is ammonium that is decoupled from 
oxidation reduction reactions so that it does not get reduced to N2.  

\begin{itemize}
\item Make sure that all of the solutes show up in the simulation results (no keypunch errors).
\item Are there solids that could precipitate based on thermodynamics ?
\end{itemize}

\begin{itemize}
\item Calculate the composition on exchange sites such that the exchanger 
is in equilibrium with the groundwater composition.  
\end{itemize}

In other words, set up conditions such that when the groundwater composition 
is reacted with the exchanger there is not net change in solution composition 
the composition on the exchange sites.  
For this step we will use thermodynamic data for ion exchange 
that are built in to the PHREEQC databases.  
This model is a reasonable model for the ion exchange assemblage in typical soils.

\begin{itemize}
\item To enter system-specific information for ion exchange 
problems use the keyword {\color{red}\textbf{EXCHANGE n}}.  
Note that, like other keywords used to enter system-specific information, 
{\color{red}\textbf{EXCHANGE}} takes an index. Consult the manual to make sure the syntax is correct.  
For 1 liter of groundwater the aquifer has 6.5 $\times$ 10$^{-3}$ moles of exchange sites.  
Enter that information and use the {\color{red}\textbf{-equilibrate solution n}} command to 
force the exchanger to be in equilibrium with the groundwater composition.
\end{itemize}

\begin{quote}
{\color{red}EXCHANGE 1 \\
\# 1	2 \\
X	6.5E-3 \\
-equilibrate with Solution n  \\
\# 1: name of the exchange site. In this case, the name in the Amm.dat database. \\
\# 2: Moles exchange sites for this exercise. }\\
\end{quote}

The command {\color{red}\textbf{-equilibrate with Solution n}} 
forces the exchanger to be in equilibrium with groundwater composition.  
Be sure to  replace {\color{red}\textbf{n}} with the index you assigned 
to the Solution with the groundwater composition.  

\begin{itemize}
\item Run the simulation.
\item Make sure there are no major changes in the redox speciation.
In other words, the redox-reactive species ($Fe(2)$, $Mn(2)$, $S(6)$, $C(4)$) 
should not have undergone oxidation or reduction to any significant extent.  
\item The dominant cations on a charge basis in solution 
are $Na$ (580 $\mu$eq/L) > $Fe(2)$ (350 $\mu$eq/L) 
  > $Ca$ (340 $\mu$eq/L) > $Mg$ (106 $\mu$eq/L) > $K$ (75 $\mu$eq/L). 
\item What is the order of dominance of cations on the exchange sites?
\item How does the concentration of $Ca$, $Fe$, and $Na$ on the exchange 
sites compare to the concentration of the $Ca$, $Fe$, and $Na$ in solution ?  
The total concentration of $Ca$, $Fe$, and $Na$ in the aquifer equals their 
concentration in groundwater plus their concentration in aqueous solution.  
What fraction of the total concentration of $Ca$, $Fe$, and $Na$ is in aqueous solution ?
\end{itemize}

Now let's react groundwater spiked with Cs with the exchange 
composition in equilibrium with groundwater.  
This simulates a batch experiment were sediments equilibrated 
with groundwater are removed and resuspended in groundwater spiked with $Cs$.

\begin{itemize}
\item First, set up a {\color{red}\textbf{SOLUTION}} 
with the composition of groundwater spiked with $Cs$.
\end{itemize}

Second, you will need to use {\color{red}\textbf{SAVE}} to 
save the composition of the exchange assemblage before the 
{\color{red}\textbf{END}} statement from the first simulation.  
After the {\color{red}\textbf{END}} statement from the first simulation, 
you will need to {\color{red}\textbf{USE}} the saved 
composition on the exchange assemblage. 
Assuming the {\color{red}\textbf{EXCHANGE}} keyword was assigned in index of 1:

\begin{quote}
{\color{red}SAVE EXCHANGE 2 \\
END \\
USE EXCHANGE 2 }\\
\end{quote}

\begin{itemize}
\item Once you have entered the composition of the groundwater spiked with Cs, run the simulation.
\end{itemize}

Note that the composition on the exchanger has not changed ! 
This is because there is no logK for $Cs$ exchange in the database.  
Assume that the selectivity of $Cs+$ is the same as K+.  
Add the necessary thermodynamic information for $Cs+$ exchange just below the 
{\color{red}\textbf{SOLUTION\_SPECIES}} entry for $Cs+$. 
That is done using the {\color{red}\textbf{EXCHANGE\_SPECIES}} keyword.  
Add the information for $Cs+$ exchange, being sure to consult the manual to get the syntax right:


\begin{quote}
{\color{red}EXCHANGE\_SPECIES \\
Cs+  +  X-  =  CsX \\
log\_k	0.7 }\\
\end{quote}

\begin{itemize}
\item Run the simulation. Make sure there is $Cs$ in the exchange assemblage.  

\item What is the concentration of $CsX$ ?  
How does it compare with the concentration of $Cs+$ remaining in solution?

\item Save the new exchange composition (with $Cs$). 
Use that exchange composition and react it again with groundwater spiked with $Cs$.  
What is the $CsX$ concentration now? 

\item Save the new exchange composition

\end{itemize}

\begin{quote}
{\color{red}SAVE EXCHANGE 2 \\
END }\\
\end{quote}

Use the exchange composition you just generated and react it again with $Cs$-spike groundwater.  
How does the $CsX$ concentration (and aqueous $Cs$ concentration) change ?  
Does it cause changes in the major ion distribution on exchange sites:

\begin{quote}
{\color{red}USE EXCHANGE 2 \\
USE Solution 2 \\
SAVE Exchange 2 \\
END }\\
\end{quote}

\begin{itemize}
\item After three equilibrations with Cs-spiked groundwater, 
is the Cs concentration on the exchanger changing much ?

\item Now, having equilibrated the exchanger three times with Cs-spike groundwater, 
run a set of three simulations where you equilibrate the exchanger with the original groundwater.  
How does the Cs concentration on the excange change ?  
If you called the original groundwater {\color{red}\textbf{SOLUTION 1}} 
and have not over-written it with another composition, you can still use it:
\end{itemize}

\begin{quote}
{\color{red}USE EXCHANGE 2 \\
\# Recalls exchange composition after \\
\# third equilibration with Cs-groundwater. \\
USE SOLUTION 1 \\
\# Original groundwater composition. \\
SAVE EXCHANGE 2 \\
END }\\
\end{quote}



